\section{Introduction}\label{sec:intro}

Instruction tuning turns a next-token predictor into an assistant, yet most of what it changes is cosmetic. When we line up a base model and its instruction-tuned sibling on the same chat-formatted text, their top-1 predictions agree on 83--85\% of response tokens, and the mean per-token \klt{} between them is only 0.21--0.24 nats. The Superficial Alignment Hypothesis~\citep{zhou2023lima} and analyses of the post-tuning distribution shift~\citep{lin2023urial} argue that alignment mostly reweights a thin band of stylistic and formatting tokens that the base model already knew how to produce. If that were the whole story, the gap between base and instruct would be almost entirely a matter of surface manners.

But recent work pushes back: instruction tuning also reshapes substantive behavior, and the ``superficial'' framing undersells what changes~\citep{revisiting2024superficial}. This leaves a sharp empirical question that prior work has not answered at the resolution of individual tokens. {\bf Where does alignment do something a simple model of the base could not have guessed?} If a divergence is predictable from cheap base-side context---this token is low-confidence, that one is high-entropy---then it is the boring, stylistic part of alignment doing its job. The interesting part is the divergence that remains \emph{unpredictable}: the positions where the instruct model moves probability mass in a way the base context gives no warning of. We argue that is precisely where alignment surprises, and that reading those surprises tells us what instruction tuning actually installs.

\para{Our approach.} We make this precise by turning ``surprise'' into a residual. For each instruction we greedily decode an \(\leq 80\)-token response with the instruct model, then teacher-force the identical token sequence (chat prompt plus response) through \emph{both} the base and the instruct model under the same chat-formatted context. This isolates the weight difference alone: same tokens, same context, only the parameters differ. At each response position we record the target \klt{} over the full vocabulary together with a handful of base-side features---position, base entropy, base top-1 probability, the base log-probability of the emitted token, token-class flags, and token length. We train a gradient-boosted tree regressor (\gbt) to predict \klt{} from those features on \dolly{}~\citep{dolly2023}, then study the \emph{residuals}---actual minus predicted \klt{}---on the disjoint \justeval{} set~\citep{lin2023urial}, 1000 instructions of which 200 carry safety tags. We repeat the whole pipeline across three scales of the \qwen{} family~\citep{qwen2025}: \qwensmall{}, \qwenmid{}, and \qwenbig{}, base and instruct.

\para{What we find.} Per-token divergence is largely predictable. The \gbt{} reaches cross-corpus \(R^2 \approx 0.60\)--\(0.68\) (0.599 / 0.680 / 0.626 across the three pairs), with Spearman \(\rho\) of 0.82--0.89 between predicted and actual \klt{}. Two features carry almost all the signal---base token log-probability and base entropy---while a position-only baseline is essentially useless (\(R^2 \leq 0.05\)) and every stylistic or lexical token flag contributes under 0.03 importance. So most of the base\(\to\)instruct shift really is anticipatable from how confident the base model already was. The residual, however, is not noise: it is structured and alignment-specific. Safety-tagged tokens carry a significantly positive mean residual in all three pairs (\eg{} \qwensmall{}: \(+0.138\) vs.\ \(-0.012\) for regular tokens), and the largest residuals concentrate on safety language, discourse and structural pivots, and persona/identity tokens---not the generic stylistic tokens that dominate the predictable shift. A telling case: after a refusal, the instruct model emits ``If you have any other questions\ldots'', but the base ranks ``If'' only third, yielding \(\klt{}\approx 19.9\)---a refusal-recovery pivot the base context never signals. Identity tokens like ``Alibaba'' (Qwen's maker) and knowledge-cutoff framing like ``October'' surface the same way.

This reframes the superficial-alignment debate quantitatively: alignment is \emph{mostly} predictable reweighting, but the unpredictable remainder is exactly the safety, framing, and persona behavior that motivates alignment in the first place.

\para{Contributions.}
\begin{itemize}[leftmargin=*,itemsep=0pt,topsep=0pt]
  \item We operationalize per-token base\(\leftrightarrow\)instruct divergence as a teacher-forced \klt{} under a shared chat context, isolating the weight difference from any decoding or prompting confound.
  \item We train a base-side \gbt{} predictor of this divergence and show it captures most of the shift (cross-corpus \(R^2 \approx 0.60\)--\(0.68\), Spearman \(\rho\) up to 0.89), driven almost entirely by base confidence and entropy rather than position or lexical flags.
  \item We define and characterize the \emph{residual}---the unpredictable divergence---and demonstrate it is structured rather than noise (Kruskal--Wallis across token classes \(p<10^{-22}\) in every pair).
  \item We show across three model scales that the residual is alignment-specific: it concentrates on safety (significant in all 3 pairs), discourse/structural pivots, and persona/identity tokens, with the safety effect shrinking as scale grows.
\end{itemize}

We review related work on the alignment delta and distribution shift in \secref{sec:related}, detail our teacher-forcing and prediction setup in \secref{sec:method}, report predictability and residual analyses in \secref{sec:results}, and discuss implications in \secref{sec:discussion}.
